\subsection{Laplace Transform of a Power Function}

According to the Laplace transform table, what is the transform of $f(t) = t^2$?

\begin{choices}
    \choice 1/s^2
    \choice 1/s^3
    \choice 2/s^2
    \correctchoice 2/s^3
\end{choices}

\begin{solution}

\subsection*{Step 1: Identify the Standard Formula}
According to the standard Laplace transform table, the transform of a power function $t^n$ for a non-negative integer $n$ is given by: $\mathcal{L}\{t^n\} = \frac{n!}{s^{n+1}}$.

\subsection*{Step 2: Substitute the Given Value}
Here, the given function is $f(t) = t^2$, which means $n = 2$. Substituting $n = 2$ into the formula gives: \mathcal{L}\{t^2\} = \frac{2!}{s^{2+1}}.

\subsection*{Step 3: Simplify the Expression}
Calculate the factorial and the exponent: $2! = 2 \times 1 = 2$ and $2 + 1 = 3$. Thus, the Laplace transform is: $\mathcal{L}\{t^2\} = \frac{2}{s^3}$.

Correct Answer: D
\end{solution}

\clearpage

\subsection{Laplace Transform with Exponential Shifting}

What is the Laplace transform of $f(t) = e^{-3t} \sin(2t)$?

\begin{choices}
    \correctchoice 2 / ((s+3)^2 + 4)
    \choice (s+3) / ((s+3)^2 + 4)
    \choice 2 / ((s-3)^2 + 4)
    \choice (s-3) / ((s-3)^2 + 4)
\end{choices}

\begin{solution}

\subsection*{Step 1: Identify the Basic Transform}
First, find the Laplace transform of the trigonometric component, $\sin(2t)$, using the standard formula $\mathcal{L}\{\sin(\omega t)\} = \frac{\omega}{s^2 + \omega^2}$: $\mathcal{L}\{\sin(2t)\} = \frac{2}{s^2 + 2^2} = \frac{2}{s^2 + 4}$.

\subsection*{Step 2: Apply the First Shifting Theorem}
The First Shifting Theorem states that if $\mathcal{L}\{f(t)\} = F(s)$, then: $\mathcal{L}\{e^{at}f(t)\} = F(s - a)$.

\subsection*{Step 3: Substitute the Value of a}
In this problem, the exponential term is $e^{-3t}$, so $a = -3$. This means we shift $s$ by replacing every $s$ with $(s - (-3)) = (s + 3)$: $\mathcal{L}\{e^{-3t} \sin(2t)\} = \left. \frac{2}{s^2 + 4} \right|_{s \to s+3} = \frac{2}{(s+3)^2 + 4}$.

Correct Answer: A
\end{solution}

\clearpage

\subsection{Inverse Laplace Transform}

Determine the inverse Laplace transform of $F(s) = \frac{1}{s-5}$.

\begin{choices}
    \choice e^{-5t}
    \correctchoice e^{5t}
    \choice 5e^t
    \choice \sin(5t)
\end{choices}

\begin{solution}

\subsection*{Step 1: Reference the Standard Inverse Formula}
The standard formula for the inverse Laplace transform of a simple linear shifted fraction is: $\mathcal{L}^{-1}\left\{\frac{1}{s-a}\right\} = e^{at}$.

\subsection*{Step 2: Match the Given Expression}
The given function is: $F(s) = \frac{1}{s-5}$. By comparing this expression with the standard template $\frac{1}{s-a}$, we find that: $a = 5$.

\subsection*{Step 3: Write out the Solution}
Substituting $a = 5$ into the exponential formula gives: $f(t) = e^{5t}$.

Correct Answer: B
\end{solution}

\clearpage

\subsection{Laplace Transform of a Derivative}

What is the Laplace transform of the derivative $f'(t)$, given the transform of $f(t)$ is $F(s)$?

\begin{choices}
    \choice F(s) / s
    \correctchoice sF(s) - f(0)
    \choice s^2F(s) - sf(0) - f'(0)
    \choice F(s) - f(0)
\end{choices}

\begin{solution}

\subsection*{Step 1: Recall the Definition/Property of Derivative Transforms}
The Laplace transform of a derivative is derived using integration by parts from the basic definition of the Laplace transform: $\mathcal{L}\{f'(t)\} = \int_{0}^{\infty} e^{-st} f'(t) \, dt$.

\subsection*{Step 2: State the Established Identity Rule}
Applying the integration by parts rule yields the standard property for the first derivative of a function: $\mathcal{L}\{f'(t)\} = sF(s) - f(0)$, where $F(s) = \mathcal{L}\{f(t)\}$ and $f(0)$ represents the initial condition of the function at $t = 0$.

Correct Answer: B
\end{solution}

\clearpage

\subsection{Definition of Laplace}

The Laplace Transform physically converts a function of time $f(t)$ into a function of the complex frequency variable $s$. It is mathematically defined strictly by the improper integral:

\begin{choices}
    \correctchoice $L\{f(t)\} = \int_{0}^{\infty} f(t) e^{st} dt$
    \choice $L\{f(t)\} = \int_{-\infty}^{\infty} f(t) e^{-st} dt$
    \choice $L\{f(t)\} = \int_{0}^{\infty} f(t) e^{-st} dt$
    \choice $L\{f(t)\} = \frac{d}{dt} [f(t) e^{-st}]$
\end{choices}

\begin{solution}

\subsection*{Step 1: Wait, let's look closer}
Actually, Option C is the standard correct definition. Let's fix the labeling. The correct integral strictly requires $e^{-st}$.

\subsection*{Step 2: The Integral}
The definition is $L\{f(t)\} = \int_0^\infty f(t) e^{-st} dt$. It takes a time-domain differential equation and converts it into a simple algebra problem in the s-domain.

Correct Answer: A
\end{solution}

\clearpage

\subsection{Principal Stresses}

A plane stress element is subjected to normal stresses $\sigma_x = 120 \text{ MPa}$ and $\sigma_y = 40 \text{ MPa}$, and a shear stress $\tau_{xy} = 30 \text{ MPa}$. Using Mohr's circle, what are the principal stresses $\sigma_1$ and $\sigma_2$?

\begin{choices}
    \choice $110 \text{ MPa}, 50 \text{ MPa}$
    \correctchoice $130 \text{ MPa}, 30 \text{ MPa}$
    \choice $150 \text{ MPa}, 10 \text{ MPa}$
    \choice $140 \text{ MPa}, 20 \text{ MPa}$
\end{choices}

\begin{solution}

\subsection*{Step 1: Find the Average Stress}
The center of Mohr's circle is $\sigma_{avg} = \frac{\sigma_x + \sigma_y}{2} = \frac{120 + 40}{2} = 80 \text{ MPa}$.

\subsection*{Step 2: Find the Radius}
$R = \sqrt{(\frac{\sigma_x - \sigma_y}{2})^2 + \tau_{xy}^2} = \sqrt{40^2 + 30^2} = 50 \text{ MPa}$.

\subsection*{Step 3: Calculate Principal Stresses}
$\sigma_1 = \sigma_{avg} + R = 80 + 50 = 130 \text{ MPa}$. $\sigma_2 = \sigma_{avg} - R = 80 - 50 = 30 \text{ MPa}$.

Correct Answer: B
\end{solution}

\clearpage

\subsection{Maximum In-Plane Shear Stress}

For a state of plane stress where $\sigma_{x} = 60\text{ MPa}$, $\sigma_{y} = -20\text{ MPa}$, and $\tau_{xy} = 30\text{ MPa}$, determine the maximum in-plane shear stress.

\begin{choices}
    \correctchoice 50 MPa
    \choice 40 MPa
    \choice 30 MPa
    \choice 70 MPa
\end{choices}

\begin{solution}

\subsection*{Step 1: Apply Radius Formula}
The maximum in-plane shear stress equals the radius $R$ of Mohr's circle: $\tau_{max} = \sqrt{(\frac{\sigma_x - \sigma_y}{2})^2 + \tau_{xy}^2}$.

\subsection*{Step 2: Substitute Values}
$\tau_{max} = \sqrt{(\frac{60 - (-20)}{2})^2 + 30^2} = \sqrt{40^2 + 30^2} = 50 \text{ MPa}$.

Correct Answer: A
\end{solution}

\clearpage

\subsection{Principal Plane Angle Calculation}

A point is subjected to $\sigma_{x} = 120\text{ MPa}$, $\sigma_{y} = 40\text{ MPa}$, and $\tau_{xy} = 30\text{ MPa}$. At what angle $\theta_p$ from the x-axis is the nearest principal plane located?

\begin{choices}
    \choice $53.1^{\circ}$
    \correctchoice $18.4^{\circ}$
    \choice $36.9^{\circ}$
    \choice $71.6^{\circ}$
\end{choices}

\begin{solution}

\subsection*{Step 1: Apply Orientation Formula}
$\tan(2\theta_p) = \frac{2\tau_{xy}}{\sigma_x - \sigma_y} = \frac{2(30)}{120 - 40} = \frac{60}{80} = 0.75$.

\subsection*{Step 2: Solve for Angle}
$2\theta_p = \arctan(0.75) \approx 36.87^{\circ} \implies \theta_p \approx 18.4^{\circ}$.

Correct Answer: B
\end{solution}

\clearpage

\subsection{Mohr's Circle Parameters}

A point in a structural member is subjected to a tensile stress of $60$ MPa in the y-direction and a compressive stress of $40$ MPa in the x-direction, with no shear stress. What are the coordinates of the center $(C, 0)$ and the radius $R$ of Mohr's circle for this state of stress?

\begin{choices}
    \choice $C = (50, 0), R = 10$
    \choice $C = (20, 0), R = 40$
    \choice $C = (10, 0), R = 100$
    \correctchoice $C = (10, 0), R = 50$
\end{choices}

\begin{solution}

\subsection*{Step 1: Identify Stress Components}
$\sigma_x = -40$ MPa (compressive), $\sigma_y = 60$ MPa (tensile), $\tau_{xy} = 0$.

\subsection*{Step 2: Calculate Center and Radius}
$C = \frac{-40 + 60}{2} = 10$. $R = \sqrt{(\frac{-40 - 60}{2})^2 + 0^2} = \sqrt{(-50)^2} = 50$.

Correct Answer: D
\end{solution}

\clearpage

\subsection{Maximum Shear Stress from Principal Stresses}

If the principal stresses at a point are σ₁ = 100 MPa and σ₂ = 40 MPa, what is the maximum in-plane shear stress?

\begin{choices}
    \correctchoice 30 MPa
    \choice 70 MPa
    \choice 60 MPa
    \choice 140 MPa
\end{choices}

\begin{solution}

\subsection*{Step 1: Max Shear Stress}
$\tau_{max} = \dfrac{\sigma_1 - \sigma_2}{2} = \dfrac{100 - 40}{2} = 30$ MPa.

Correct Answer: A
\end{solution}

\clearpage

\subsection{Peritectic Reaction in Binary Alloys}

Which of the following describes a peritectic reaction in a binary alloy system during cooling?

\begin{choices}
    \choice A liquid phase transforms into two different solid phases.
    \correctchoice A liquid phase and a solid phase transform into a different solid phase.
    \choice A solid phase transforms into two different solid phases.
    \choice Two solid phases transform into a single, different solid phase.
\end{choices}

\begin{solution}

\subsection*{Step 1: Analyzing the Peritectic Reaction}
A peritectic reaction occurs when a liquid phase ($L$) and a solid phase (e.g., $\alpha$) react at a specific temperature to form a new, different solid phase (e.g., $\beta$). This is represented as: $L + \alpha \xrightarrow{\text{cooling}} \beta$.

\subsection*{Step 2: Comparison with Other Reactions}
Eutectic: $L \rightarrow \alpha + \beta$; Eutectoid: $\gamma \rightarrow \alpha + \beta$; Peritectoid: $\alpha + \beta \rightarrow \gamma$.

Correct Answer: B
\end{solution}

\clearpage

\subsection{Invariant Points in a Phase Diagram}

In the iron-iron carbide ($Fe - Fe_3C$) phase diagram, the eutectoid reaction occurs at a specific temperature and carbon concentration. What are the correct values for this invariant point?

\begin{choices}
    \choice $1147^\circ\text{C}$ and $4.30 \text{ wt\% C}$
    \choice $912^\circ\text{C}$ and $0.022 \text{ wt\% C}$
    \choice $1495^\circ\text{C}$ and $0.16 \text{ wt\% C}$
    \correctchoice $727^\circ\text{C}$ and $0.76 \text{ wt\% C}$
\end{choices}

\begin{solution}

\subsection*{Step 1: Eutectoid Reaction Definition}
The eutectoid reaction in the $Fe-Fe_3C$ system involves solid Austenite ($\gamma$) transforming into Ferrite ($\alpha$) and Cementite ($Fe_3C$) at $727^\circ\text{C}$ and $0.76 \text{ wt\% C}$.

\subsection*{Step 2: Comparison of Invariant Points}
Eutectic point occurs at $1147^\circ\text{C}$ and $4.30 \text{ wt\% C}$. Peritectic point occurs at $1495^\circ\text{C}$ and $0.16 \text{ wt\% C}$.

Correct Answer: D
\end{solution}

\clearpage

\subsection{Lever Rule}

In a binary eutectic phase diagram for components $A$ and $B$, the eutectic point occurs at $60\% B$ and $727^\circ\text{C}$. A $40\% B$ alloy is cooled from the liquid phase to just below the eutectic temperature. What is the mass fraction of the primary $\alpha$ phase (where the maximum solubility of $B$ in $A$ is $10\%$ at the eutectic temperature)?

\begin{choices}
    \correctchoice 0.40
    \choice 0.67
    \choice 0.60
    \choice 0.25
\end{choices}

\begin{solution}

\subsection*{Step 1: Identify the Compositions}
$C_0$ (Alloy) = $40\% B$, $C_{\alpha}$ (Solubility limit) = $10\% B$, $C_E$ (Eutectic/Liquid phase) = $60\% B$.

\subsection*{Step 2: Apply the Lever Rule}
The mass fraction of primary $\alpha$ ($W_{\alpha'}$) just above the eutectic temperature is: $W_{\alpha'} = \dfrac{C_E - C_0}{C_E - C_{\alpha}}$.

\subsection*{Step 3: Perform Calculation}
$W_{\alpha'} = \dfrac{60 - 40}{60 - 10} = \dfrac{20}{50} = 0.40$.

Correct Answer: A
\end{solution}

\clearpage

\subsection{Lever Rule — Phase Fraction}

In a binary alloy at a temperature where α and liquid phases coexist, if the alloy composition is 40% B, the α boundary is at 20% B, and the liquid boundary is at 60% B, what fraction of the alloy is liquid?

\begin{choices}
    \correctchoice 0.50
    \choice 0.25
    \choice 0.75
    \choice 0.40
\end{choices}

\begin{solution}

\subsection*{Step 1: Lever Rule}
Fraction liquid = $\dfrac{C_0 - C_\alpha}{C_L - C_\alpha} = \dfrac{40 - 20}{60 - 20} = \dfrac{20}{40} = 0.50$.

Correct Answer: A
\end{solution}

\clearpage

\subsection{Transformer Turns Ratio}

A transformer has 500 primary turns and 100 secondary turns. If the primary voltage is 240V, what is the secondary voltage?

\begin{choices}
    \correctchoice 48 V
    \choice 1200 V
    \choice 240 V
    \choice 24 V
\end{choices}

\begin{solution}

\subsection*{Step 1: Turns Ratio}
$\dfrac{V_s}{V_p} = \dfrac{N_s}{N_p} \Rightarrow V_s = 240 \times \dfrac{100}{500} = 48$ V.

Correct Answer: A
\end{solution}

\clearpage

\subsection{Cycle Length}

In the timing design of a standardized traffic signal intersection, the 'Cycle Length' ($C$) is mathematically and physically defined as the:

\begin{choices}
    \choice Total time exactly required for one vehicle to cross the entire intersection
    \choice Time duration of the yellow clearance interval
    \choice Time taken for a pedestrian to walk across the crosswalk
    \correctchoice Total time required for one complete sequence of all signal phases (e.g., from the start of Main Street Green to the next start of Main Street Green)
\end{choices}

\begin{solution}

\subsection*{Step 1: Phasing}
A signal rotates through phases (North/South Green, N/S Yellow, N/S Red, East/West Green, E/W Yellow, E/W Red).

\subsection*{Step 2: The Cycle}
The Cycle Length is the total time, in seconds, it takes the computer to complete that entire loop before repeating. A typical cycle length is between 60 and 120 seconds.

Correct Answer: D
\end{solution}

\clearpage

\subsection{Ideal Transformer Voltage}

For an ideal transformer, $V_1/V_2 = $:

\begin{choices}
    \choice $I_1/I_2$
    \choice $N_2/N_1$
    \correctchoice $N_1/N_2$
    \choice $(N_1/N_2)^2$
\end{choices}

\begin{solution}

\subsection*{Step 1: Ratio}
Voltage ratio is proportional to turns ratio.

Correct Answer: C
\end{solution}

\clearpage

\subsection{Ideal Transformer Turn Ratio Voltage Scaling}

An ideal transformer has $N_1 = 500$ turns on the primary winding and $N_2 = 100$ turns on the secondary winding. If a voltage of $V_1 = 120\text{ V}$ AC is applied to the primary, calculate the secondary open-circuit voltage ($V_2$).

\begin{choices}
    \correctchoice 24.0 V
    \choice 600.0 V
    \choice 12.0 V
    \choice 48.0 V
\end{choices}

\begin{solution}

\subsection*{Step 1: Identify Transformer Voltage Relation}
For an ideal transformer: $\frac{V_1}{V_2} = \frac{N_1}{N_2}$.

\subsection*{Step 2: Solve for V2}
$V_2 = V_1 \left(\frac{N_2}{N_1}\right) = 120 \times \left(\frac{100}{500}\right) = 120 \times 0.20 = 24.0\text{ V}$.

Correct Answer: A
\end{solution}

\clearpage

\subsection{Ideal Transformer Current Scaling}

For the transformer in the previous question ($N_1/N_2 = 5$), if a load connected to the secondary draws $I_2 = 10\text{ A}$ of current, what is the current $I_1$ drawn from the primary supply source?

\begin{choices}
    \correctchoice 2.0 A
    \choice 50.0 A
    \choice 10.0 A
    \choice 5.0 A
\end{choices}

\begin{solution}

\subsection*{Step 1: Identify Transformer Current Relation}
For an ideal transformer, input power equals output power: $V_1 I_1 = V_2 I_2 \implies \frac{I_1}{I_2} = \frac{N_2}{N_1}$.

\subsection*{Step 2: Solve for I1}
$I_1 = I_2 \left(\frac{N_2}{N_1}\right) = 10 \times \left(\frac{100}{500}\right) = 2.0\text{ A}$.

Correct Answer: A
\end{solution}

\clearpage

\subsection{Transformer Impedance Reflection Sizing}

A secondary load impedance is $\mathbf{Z}_L = 4\text{ \Omega}$ connected to a transformer with turn ratio $a = N_1/N_2 = 10$. Calculate the reflected impedance $\mathbf{Z}'_L$ seen looking into the primary terminals.

\begin{choices}
    \correctchoice 400 \Omega
    \choice 40 \Omega
    \choice 0.04 \Omega
    \choice 200 \Omega
\end{choices}

\begin{solution}

\subsection*{Step 1: Identify Impedance Reflection Formula}
The impedance reflected to the primary is: $\mathbf{Z}'_L = a^2 \mathbf{Z}_L$, where $a = \frac{N_1}{N_2}$.

\subsection*{Step 2: Calculate Reflected Impedance}
$\mathbf{Z}'_L = 10^2 \times 4\text{ \Omega} = 100 \times 4 = 400\text{ \Omega}$.

Correct Answer: A
\end{solution}

\clearpage

\subsection{Transformer Dot Convention Polarities}

The dot convention is used to determine the relative polarities of mutually coupled coils. If current enters the dotted terminal of the primary winding, what is the polarity of the induced voltage at the dotted terminal of the secondary winding?

\begin{choices}
    \correctchoice Positive relative to its non-dotted terminal
    \choice Negative relative to its non-dotted terminal
    \choice Zero throughout
    \choice Inverted at exactly 180 degrees phase shift
\end{choices}

\begin{solution}

\subsection*{Step 1: Apply Dot Convention Rules}
The dot convention states that if a current enters the dotted terminal of one coil, it induces a positive voltage polarity at the dotted terminal of the second mutually coupled coil.

Correct Answer: A
\end{solution}

\clearpage

\subsection{Autotransformer Power Rating Sizing}

An autotransformer is connected to step down a voltage from $240\text{ V}$ to $200\text{ V}$. If the secondary load draws $5\text{ kVA}$, what portion of this power is transferred by direct electrical conduction (conductive power) rather than magnetic induction?

\begin{choices}
    \correctchoice 4.17 kVA
    \choice 0.83 kVA
    \choice 2.50 kVA
    \choice 5.00 kVA
\end{choices}

\begin{solution}

\subsection*{Step 1: Identify Power Transfer Formula}
For an autotransformer, the power transferred by induction is: $S_{ind} = S_{total} \left(1 - \frac{V_{low}}{V_{high}}\right)$.

\subsection*{Step 2: Calculate Inductive Power}
$S_{ind} = 5\text{ kVA} \times \left(1 - \frac{200}{240}\right) = 5 \times \left(1 - 0.833\right) = 5 \times 0.1667 = 0.833\text{ kVA}$.

\subsection*{Step 3: Calculate Conductive Power}
$S_{cond} = S_{total} - S_{ind} = 5.0 - 0.833 = 4.167\text{ kVA}$.

Correct Answer: A
\end{solution}

\clearpage

\subsection{Real Transformer Equivalent Core Loss Resistance}

In the equivalent circuit model of a real transformer, which parameter represents the eddy current and hysteresis core losses?

\begin{choices}
    \correctchoice Shunt resistance (R_c) in parallel with magnetizing reactance (X_m)
    \choice Series winding resistance (R_eq)
    \choice Series leakage reactance (X_eq)
    \choice The ideal transformation turn ratio directly
\end{choices}

\begin{solution}

\subsection*{Step 1: Analyze Real Transformer Shunt Branch}
Core losses (hysteresis and eddy currents) depend on flux density and are modeled by a shunt resistor $R_c$ connected across the primary induced voltage. The magnetizing current that establishes flux is modeled by parallel shunt reactance $X_m$.

Correct Answer: A
\end{solution}

\clearpage

\subsection{Three-phase Transformer Delta-Wye Voltage Ratio}

A three-phase transformer bank is connected in Delta-Wye ($\Delta-Y$) with a turn ratio per phase of $a = N_{primary}/N_{secondary} = 5$. If the primary line-to-line voltage is $4160\text{ V}$, what is the secondary line-to-line voltage?

\begin{choices}
    \correctchoice 1441 V
    \choice 832 V
    \choice 2400 V
    \choice 480 V
\end{choices}

\begin{solution}

\subsection*{Step 1: Analyze Delta-Wye Phase Relations}
Primary is Delta, so phase voltage equals line voltage: $V_{\phi,p} = 4160\text{ V}$.
Secondary phase voltage: $V_{\phi,s} = \frac{V_{\phi,p}}{a} = \frac{4160}{5} = 832\text{ V}$.

\subsection*{Step 2: Calculate Line-to-Line Secondary Voltage}
Secondary is Wye, so line-to-line voltage is: $V_{L,s} = \sqrt{3} V_{\phi,s} = \sqrt{3} \times 832 = 1441\text{ V}$.

Correct Answer: A
\end{solution}

\clearpage

\subsection{Real Transformer Efficiency}

A $10\text{ kVA}$ transformer has core losses of $150\text{ W}$ and full-load copper losses of $250\text{ W}$. Calculate the efficiency of the transformer operating at full load and unity power factor.

\begin{choices}
    \correctchoice 96.15%
    \choice 97.56%
    \choice 98.20%
    \choice 94.80%
\end{choices}

\begin{solution}

\subsection*{Step 1: Calculate Output Power}
$P_{out} = S \cdot pf = 10000 \times 1.0 = 10000\text{ W}$.

\subsection*{Step 2: Calculate Total Losses}
$P_{losses} = P_{core} + P_{copper} = 150 + 250 = 400\text{ W}$.

\subsection*{Step 3: Calculate Efficiency}
$\eta = \frac{P_{out}}{P_{out} + P_{losses}} = \frac{10000}{10000 + 400} = \frac{10000}{10400} = 96.15\%$.

Correct Answer: A
\end{solution}

\clearpage

\subsection{Ideal Transformer Apparent Impedance Matching}

A transformer is used for impedance matching between an audio amplifier with output resistance $R_s = 800\text{ \Omega}$ and a speaker with resistance $R_L = 8\text{ \Omega}$. What turn ratio $a = N_1/N_2$ is required to match these impedances?

\begin{choices}
    \correctchoice 10.0
    \choice 100.0
    \choice 5.0
    \choice 20.0
\end{choices}

\begin{solution}

\subsection*{Step 1: Identify Impedance Matching Condition}
For perfect matching, reflected impedance must equal source resistance: $R_s = a^2 R_L$.

\subsection*{Step 2: Solve for Turn Ratio}
$a^2 = \frac{R_s}{R_L} = \frac{800}{8} = 100 \implies a = \sqrt{100} = 10.0$.

Correct Answer: A
\end{solution}

\clearpage

\subsection{Step Function}

The Laplace transform of the unit step function $u(t)$ is:

\begin{choices}
    \choice 1
    \choice $1/s^2$
    \choice $s$
    \correctchoice $1/s$
\end{choices}

\begin{solution}

\subsection*{Step 1: Formula}
Basic Laplace transform pair.

Correct Answer: D
\end{solution}

\clearpage

\subsection{Impulse Function}

The Laplace transform of the Dirac delta function $delta(t)$ is:

\begin{choices}
    \choice $1/s$
    \choice $s$
    \choice $infty$
    \correctchoice 1
\end{choices}

\begin{solution}

\subsection*{Step 1: Formula}
Transform of an impulse is unity.

Correct Answer: D
\end{solution}

\clearpage

\subsection{Final Value Theorem}

The Final Value Theorem $lim_{t 	o infty} f(t)$ is given by:

\begin{choices}
    \choice $lim_{s 	o infty} sF(s)$
    \choice $lim_{s 	o 0} F(s)$
    \correctchoice $lim_{s 	o 0} sF(s)$
    \choice $lim_{s 	o infty} F(s)$
\end{choices}

\begin{solution}

\subsection*{Step 1: Condition}
All poles of $sF(s)$ must be in the LHP.

Correct Answer: C
\end{solution}

\clearpage

\subsection{Integration property}

The Laplace transform of $int f(t) dt$ is:

\begin{choices}
    \correctchoice $F(s) / s$
    \choice $F(s) - f(0)$
    \choice $s F(s)$
    \choice $F'(s)$
\end{choices}

\begin{solution}

\subsection*{Step 1: Rule}
Integration in time is division by $s$ in Laplace domain.

Correct Answer: A
\end{solution}

\clearpage

\subsection{Laplace Transform of a Cosine Function}

What is the Laplace transform of the function $f(t) = e^{-3t} \cos(4t) u(t)$?

\begin{choices}
    \choice $\frac{4}{(s+3)^2 + 16}$
    \correctchoice $\frac{s+3}{(s+3)^2 + 16}$
    \choice $\frac{s+3}{(s-3)^2 + 16}$
    \choice $\frac{s}{(s+3)^2 + 16}$
\end{choices}

\begin{solution}

\subsection*{Step 1: Recall Laplace of Cosine}
$\mathcal{L}\{\cos(\omega t)\} = \frac{s}{s^2 + \omega^2}$

\subsection*{Step 2: Apply Frequency Shifting Property}
The shifting property states:
$\mathcal{L}\{e^{-at} g(t)\} = G(s+a)$

\subsection*{Step 3: Substitute Values}
Given $a = 3$ and $\omega = 4$:
$\mathcal{L}\{e^{-3t} \cos(4t)\} = \frac{s+3}{(s+3)^2 + 4^2} = \frac{s+3}{(s+3)^2 + 16}$

Correct Answer: B
\end{solution}

\clearpage

\subsection{Final Value Theorem Application}

A system's output Laplace transform is $Y(s) = \frac{5(s+2)}{s(s^2 + 3s + 2)}$. What is the final value of the output $y(\infty) = \lim_{t \to \infty} y(t)$?

\begin{choices}
    \choice 10
    \choice 0
    \correctchoice 5
    \choice Infinity
\end{choices}

\begin{solution}

\subsection*{Step 1: State Final Value Theorem}
If $s Y(s)$ has all its poles in the Left-Half Plane (LHP), then:
$y(\infty) = \lim_{s \to 0} s Y(s)$

\subsection*{Step 2: Check Pole Locations of $s Y(s)$}
$s Y(s) = \frac{5(s+2)}{s^2 + 3s + 2} = \frac{5(s+2)}{(s+1)(s+2)} = \frac{5}{s+1}$
The pole is at $s = -1$ (LHP), so the theorem is valid.

\subsection*{Step 3: Evaluate Limit}
$y(\infty) = \lim_{s \to 0} \frac{5}{s+1} = 5$

Correct Answer: C
\end{solution}

\clearpage

\subsection{Initial Value Theorem Application}

Given the Laplace transform $X(s) = \frac{2s + 5}{s^2 + 4s + 13}$, what is the initial value $x(0^+) = \lim_{t \to 0^+} x(t)$?

\begin{choices}
    \choice 5
    \choice 0
    \choice 0.38
    \correctchoice 2
\end{choices}

\begin{solution}

\subsection*{Step 1: State Initial Value Theorem}
The initial value $x(0^+)$ can be found using:
$x(0^+) = \lim_{s \to \infty} s X(s)$

\subsection*{Step 2: Evaluate Limit}
$s X(s) = \frac{s(2s + 5)}{s^2 + 4s + 13} = \frac{2s^2 + 5s}{s^2 + 4s + 13}$
Divide numerator and denominator by $s^2$:
$\lim_{s \to \infty} \frac{2 + 5/s}{1 + 4/s + 13/s^2} = \frac{2 + 0}{1 + 0 + 0} = 2$

Correct Answer: D
\end{solution}

\clearpage

\subsection{Laplace Transform of a Derivative}

If $X(s)$ is the Laplace transform of $x(t)$, what is the Laplace transform of the second derivative $\frac{d^2 x(t)}{dt^2}$ in terms of initial conditions $x(0^-)$ and $x'(0^-)$?

\begin{choices}
    \correctchoice $s^2 X(s) - s x(0^-) - x'(0^-)$
    \choice $s^2 X(s) - x(0^-) - s x'(0^-)$
    \choice $s^2 X(s)$
    \choice $\frac{X(s)}{s^2}$
\end{choices}

\begin{solution}

\subsection*{Step 1: Recall Differentiation Property}
For a first derivative:
$\mathcal{L}\{x'(t)\} = s X(s) - x(0^-)$
For a second derivative, apply the rule twice:
$\mathcal{L}\{x''(t)\} = s \mathcal{L}\{x'(t)\} - x'(0^-)$
$\mathcal{L}\{x''(t)\} = s (s X(s) - x(0^-)) - x'(0^-) = s^2 X(s) - s x(0^-) - x'(0^-)$

Correct Answer: A
\end{solution}

\clearpage

\subsection{Laplace Transform of Unit Ramp}

What is the Laplace transform of the unit ramp function $f(t) = t u(t)$?

\begin{choices}
    \choice $\frac{1}{s}$
    \correctchoice $\frac{1}{s^2}$
    \choice $s$
    \choice $\frac{2}{s^3}$
\end{choices}

\begin{solution}

\subsection*{Step 1: Calculate the Laplace Transform Integral}
$F(s) = \int_{0}^{\infty} t e^{-st}\, dt$
Using integration by parts ($u = t$, $dv = e^{-st} dt$):
$F(s) = \left[ -\frac{t}{s} e^{-st} \right]_{0}^{\infty} + \frac{1}{s} \int_{0}^{\infty} e^{-st}\, dt$
$F(s) = 0 + \frac{1}{s} \left[ -\frac{1}{s} e^{-st} \right]_{0}^{\infty} = \frac{1}{s^2}$

Correct Answer: B
\end{solution}

\clearpage

\subsection{Laplace Transform Time-Delay Property}

If the Laplace transform of $x(t) u(t)$ is $X(s)$, what is the Laplace transform of the time-delayed signal $y(t) = x(t-t_0) u(t-t_0)$ where $t_0 > 0$?

\begin{choices}
    \choice $e^{s t_0} X(s)$
    \choice $X(s - t_0)$
    \correctchoice $e^{-s t_0} X(s)$
    \choice $\frac{X(s)}{s t_0}$
\end{choices}

\begin{solution}

\subsection*{Step 1: Apply Shifting Theorem}
The time-delay theorem states that shifting a signal in time by $t_0$ (delayed) corresponds to multiplying its Laplace transform by $e^{-s t_0}$:
$\mathcal{L}\{x(t - t_0) u(t - t_0)\} = e^{-s t_0} X(s)$

Correct Answer: C
\end{solution}

\clearpage

\subsection{The S-Domain}

The Convolution Integral is a brutal, agonizing calculus nightmare to solve by hand in the Time Domain ($t$). To physically avoid this nightmare, engineers use the Laplace Transform to shift the math into the 's-Domain'. In the s-Domain, the horrific Convolution integral mathematically collapses into simple:

\begin{choices}
    \choice Subtraction
    \choice Addition
    \correctchoice Basic, 3rd-grade Algebraic MULTIPLICATION ($Y(s) = X(s) \cdot H(s)$)
    \choice Calculus derivatives
\end{choices}

\begin{solution}

\subsection*{Step 1: The Math Trick}
Trying to convolve two complex sine waves requires $5$ pages of calculus integration.

\subsection*{Step 2: The Laplace Shortcut}
You look up the Laplace translation of your input signal in a table. You multiply it by the system's Transfer Function. You then do an 'Inverse Laplace' to translate the answer back to the Time Domain. The Laplace Transform mathematically turns horrifying differential calculus into middle-school algebra.

Correct Answer: C
\end{solution}

\clearpage

\subsection{Transfer Functions (H(s))}

In the s-Domain, the 'Transfer Function' $H(s)$ is the mathematical absolute god-equation of the physical system. It is physically defined as the absolute mathematical ratio of:

\begin{choices}
    \choice Time divided by space
    \choice Voltage divided by current
    \correctchoice The Laplace Transform of the OUTPUT strictly divided by the Laplace Transform of the INPUT ($Y(s) / X(s)$), assuming absolutely zero initial energy in the system
    \choice Heat divided by friction
\end{choices}

\begin{solution}

\subsection*{Step 1: The Black Box}
If you hand an engineer a sealed black box of electronics, they don't care what resistors are inside. They only care what happens when they plug it in.

\subsection*{Step 2: The Ratio}
If they put $5\text{V}$ in, and get $10\text{V}$ out, the mathematical Transfer Function acts like a '$	imes 2$' multiplier. By analyzing $H(s)$, the engineer mathematically knows everything the box will ever physically do, without ever opening it.

Correct Answer: C
\end{solution}

\clearpage

\subsection{Low-pass Filter}

A low-pass filter passes:

\begin{choices}
    \choice Frequencies above a cutoff
    \correctchoice Frequencies below a cutoff
    \choice DC only
    \choice Only a specific range
\end{choices}

\begin{solution}

\subsection*{Step 1: Function}
Blocks high-frequency noise.

Correct Answer: B
\end{solution}

\clearpage

\subsection{Energy Signal}

An energy signal has:

\begin{choices}
    \correctchoice Finite energy and zero average power
    \choice Zero energy
    \choice Infinite power
    \choice Infinite energy and finite power
\end{choices}

\begin{solution}

\subsection*{Step 1: Definition}
Pulse-like signals are typically energy signals.

Correct Answer: A
\end{solution}

\clearpage

\subsection{Windowing}

The purpose of windowing in signal analysis is to:

\begin{choices}
    \choice Increase resolution
    \choice Amplify the signal
    \choice Filter out DC
    \correctchoice Reduce spectral leakage
\end{choices}

\begin{solution}

\subsection*{Step 1: Technique}
Tapering the edges of a finite-length signal.

Correct Answer: D
\end{solution}

\clearpage

\subsection{High-pass Filter}

A high-pass filter is used to:

\begin{choices}
    \choice Remove high-frequency noise
    \correctchoice Remove DC or low-frequency offsets
    \choice Limit bandwidth
    \choice Smooth a signal
\end{choices}

\begin{solution}

\subsection*{Step 1: Application}
Useful for edge detection or removing drift.

Correct Answer: B
\end{solution}

\clearpage

\subsection{Band-pass Filter}

A band-pass filter passes frequencies:

\begin{choices}
    \choice All frequencies
    \correctchoice Within a specific range $[f_1, f_2]$
    \choice None of the above
    \choice Outside a specific range
\end{choices}

\begin{solution}

\subsection*{Step 1: Function}
Combines low-pass and high-pass characteristics.

Correct Answer: B
\end{solution}

\clearpage

\subsection{Duality}

The Fourier transform of a rectangular pulse in time is a:

\begin{choices}
    \choice Gaussian function
    \choice Triangular pulse
    \choice Delta function
    \correctchoice Sinc function in frequency
\end{choices}

\begin{solution}

\subsection*{Step 1: Pair}
Rect $leftrightarrow$ Sinc.

Correct Answer: D
\end{solution}

\clearpage

\subsection{Filter Type: Butterworth Characteristics}

Which analog filter type is characterized by having a maximally flat magnitude response in both the passband and the stopband?

\begin{choices}
    \choice Chebyshev Type I filter
    \choice Chebyshev Type II filter
    \correctchoice Butterworth filter
    \choice Elliptic filter
\end{choices}

\begin{solution}

\subsection*{Step 1: Compare Standard Filter Profiles}
Butterworth: Maximally flat in passband/stopband, slow roll-off.
Chebyshev Type I: Ripple in passband, flat in stopband, faster roll-off.
Chebyshev Type II: Flat in passband, ripple in stopband.
Elliptic (Cauer): Ripple in both passband and stopband, steepest roll-off.

Correct Answer: C
\end{solution}

\clearpage

\subsection{Group Delay in Filters}

What is the physical meaning of the 'group delay' of a filter?

\begin{choices}
    \correctchoice The rate of change of phase shift with respect to frequency
    \choice The time taken for the filter coefficients to converge
    \choice The amplitude attenuation of a group of signals
    \choice The phase velocity at zero frequency
\end{choices}

\begin{solution}

\subsection*{Step 1: Define Group Delay}
Group delay $\tau_g(\omega)$ is the negative derivative of the phase response with respect to angular frequency:
$\tau_g(\omega) = -\frac{d\theta(\omega)}{d\omega}$
It measures the average envelope delay experienced by a narrow band of frequency components passing through the filter.

Correct Answer: A
\end{solution}

\clearpage

\subsection{Energy vs Power Signals}

A signal is classified as a 'power signal' if its total energy ($E$) and average power ($P$) satisfy which of the following?

\begin{choices}
    \choice $0 < E < \infty$ and $P = 0$
    \choice $E = \infty$ and $P = \infty$
    \correctchoice $E = \infty$ and $0 < P < \infty$
    \choice $E = 0$ and $P = 0$
\end{choices}

\begin{solution}

\subsection*{Step 1: Define Signal Energy and Power Classes}
Energy Signal: Has finite, non-zero energy ($0 < E < \infty$) which implies its average power is $P = 0$.
Power Signal: Has finite, non-zero average power ($0 < P < \infty$) which implies its total energy is $E = \infty$ (typically periodic signals).

Correct Answer: C
\end{solution}

\clearpage

\subsection{Even and Odd Parts of Signal}

Given a signal $x(t) = e^{-2t} u(t)$, what is its even part $x_e(t)$?

\begin{choices}
    \choice $e^{-2t}$
    \choice $\frac{1}{2} (e^{-2t} + e^{2t})$
    \choice $\frac{1}{2} e^{-2t} u(t)$
    \correctchoice $\frac{1}{2} e^{-2|t|}$
\end{choices}

\begin{solution}

\subsection*{Step 1: Recall Even Part Formula}
The even part $x_e(t)$ of any signal is:
$x_e(t) = \frac{x(t) + x(-t)}{2}$

\subsection*{Step 2: Substitute $x(t)$}
Given $x(t) = e^{-2t} u(t)$:
$x_e(t) = \frac{e^{-2t} u(t) + e^{2t} u(-t)}{2}$
Since $u(t)$ and $u(-t)$ are mutually exclusive except at $t=0$, this can be compactified as:
$x_e(t) = \frac{1}{2} e^{-2|t|}\quad \text{for all } t$

Correct Answer: D
\end{solution}

\clearpage

\subsection{Unit Impulse Sifting Property}

What is the value of the integral $\int_{-\infty}^{\infty} (t^3 + 5t^2 + 2) \delta(t - 2)\, dt$?

\begin{choices}
    \correctchoice 30
    \choice 2
    \choice 16
    \choice 0
\end{choices}

\begin{solution}

\subsection*{Step 1: Apply Sifting Property of Delta Function}
The sifting property states:
$\int_{-\infty}^{\infty} g(t) \delta(t - t_0)\, dt = g(t_0)$

\subsection*{Step 2: Evaluate at $t_0 = 2$}
Given $g(t) = t^3 + 5t^2 + 2$ and $t_0 = 2$:
$g(2) = 2^3 + 5(2)^2 + 2 = 8 + 20 + 2 = 30$

Correct Answer: A
\end{solution}

\clearpage

\subsection{Causal vs Anti-Causal Signals}

A continuous-time signal $x(t)$ is classified as strictly anti-causal if:

\begin{choices}
    \choice $x(t) = 0$ for all $t < 0$
    \correctchoice $x(t) = 0$ for all $t > 0$
    \choice $x(t) = x(-t)$
    \choice $x(t)$ has no real parts
\end{choices}

\begin{solution}

\subsection*{Step 1: Define Signal Causality Classes}
Causal Signal: $x(t) = 0$ for $t < 0$.
Anti-Causal Signal: $x(t) = 0$ for $t > 0$ (exists only in the past).
Non-Causal Signal: Exists in both positive and negative time.

Correct Answer: B
\end{solution}

\clearpage

\subsection{Orthogonal Signals Definition}

Two continuous-time signals $x_1(t)$ and $x_2(t)$ are said to be orthogonal over an interval $[t_1, t_2]$ if:

\begin{choices}
    \choice $\int_{t_1}^{t_2} [x_1(t) + x_2(t)]\, dt = 0$
    \choice $x_1(t) = x_2(t - t_0)$
    \correctchoice $\int_{t_1}^{t_2} x_1(t) x_2^*(t)\, dt = 0$
    \choice $x_1(t) \cdot x_2(t) = 1$
\end{choices}

\begin{solution}

\subsection*{Step 1: Define Orthogonality}
Two signals are orthogonal if their inner product over the specified interval is exactly zero. For complex signals, the inner product is:
$\langle x_1, x_2 \rangle = \int_{t_1}^{t_2} x_1(t) x_2^*(t)\, dt = 0$

Correct Answer: C
\end{solution}

\clearpage

\subsection{Determinism in Signals}

A signal whose future values can be predicted exactly with a mathematical formula is classified as:

\begin{choices}
    \choice Random (stochastic) signal
    \choice Power signal
    \choice Causal signal
    \correctchoice Deterministic signal
\end{choices}

\begin{solution}

\subsection*{Step 1: Define Determinism}
Deterministic Signal: Values are completely specified at any time instant by a formula (e.g., $x(t) = A \sin(\omega t)$).
Random Signal: Values exhibit uncertainty and must be analyzed statistically (e.g., thermal noise).

Correct Answer: D
\end{solution}

\clearpage

\subsection{Energy of Exponential Pulse}

What is the total energy $E$ of the signal $x(t) = e^{-3t} u(t)$?

\begin{choices}
    \correctchoice $1/6$
    \choice $1/3$
    \choice Infinity
    \choice $1/9$
\end{choices}

\begin{solution}

\subsection*{Step 1: Write Energy Integral}
$E = \int_{-\infty}^{\infty} |x(t)|^2\, dt = \int_{0}^{\infty} (e^{-3t})^2\, dt = \int_{0}^{\infty} e^{-6t}\, dt$

\subsection*{Step 2: Evaluate Integral}
$E = \left[ -\frac{1}{6} e^{-6t} \right]_{0}^{\infty} = 0 - \left( -\frac{1}{6} \right) = \frac{1}{6}$

Correct Answer: A
\end{solution}

\clearpage

\subsection{Signal Scaling and Time Compression}

If a signal $x(t)$ is replaced by $x(3t)$, this operation is known as:

\begin{choices}
    \choice Time expansion by a factor of 3
    \correctchoice Time compression by a factor of 3
    \choice Time delay of 3 seconds
    \choice Amplitude scaling by a factor of 3
\end{choices}

\begin{solution}

\subsection*{Step 1: Explain Time Scaling}
In time scaling $x(at)$:
- If $|a| > 1$, the signal is compressed in time by a factor of $a$.
- If $|a| < 1$, the signal is expanded/stretched in time.
Here, $a = 3$, so it is a time compression.

Correct Answer: B
\end{solution}

\clearpage

\subsection{Unit Step and Impulse Relation}

The relationship between the unit impulse function $\delta(t)$ and the unit step function $u(t)$ is represented by which of the following?

\begin{choices}
    \choice $u(t) = \frac{d\delta(t)}{dt}$
    \choice $\delta(t) = \int_{-\infty}^{t} u(\tau)\, d\tau$
    \correctchoice $\delta(t) = \frac{du(t)}{dt}$
    \choice $\delta(t) = u(t) * u(t)$
\end{choices}

\begin{solution}

\subsection*{Step 1: Relate Step and Impulse}
The unit step function $u(t)$ is the integral of the unit impulse function:
$u(t) = \int_{-\infty}^{t} \delta(\tau)\, d\tau$
By the Fundamental Theorem of Calculus, the derivative of $u(t)$ is the impulse function:
$\frac{du(t)}{dt} = \delta(t)$

Correct Answer: C
\end{solution}

\clearpage

\subsection{LTI System Frequency Response Definition}

The frequency response $H(j\omega)$ of an LTI system is the Fourier transform of its:

\begin{choices}
    \choice Step response $s(t)$
    \correctchoice Impulse response $h(t)$
    \choice Input signal $x(t)$
    \choice Transfer function $H(s)$
\end{choices}

\begin{solution}

\subsection*{Step 1: Define Frequency Response}
The frequency response $H(j\omega)$ describes the system behavior in the frequency domain. It is defined as the Fourier transform of the time-domain impulse response $h(t)$:
$H(j\omega) = \mathcal{F}\{h(t)\} = \int_{-\infty}^{\infty} h(t) e^{-j\omega t}\, dt$

Correct Answer: B
\end{solution}

\clearpage

\subsection{Phase Delay vs Group Delay}

A system has a phase response $\theta(\omega) = -3\omega$. What are the phase delay ($\tau_p$) and group delay ($\tau_g$) of this system?

\begin{choices}
    \choice $\tau_p = 3\ \text{s}$, $\tau_g = 0\ \text{s}$
    \choice $\tau_p = 0\ \text{s}$, $\tau_g = 3\ \text{s}$
    \correctchoice $\tau_p = 3\ \text{s}$, $\tau_g = 3\ \text{s}$ (Linear Phase)
    \choice $\tau_p = -3\ \text{s}$, $\tau_g = -3\ \text{s}$
\end{choices}

\begin{solution}

\subsection*{Step 1: Define Phase and Group Delays}
$\tau_p(\omega) = -\frac{\theta(\omega)}{\omega}$
$\tau_g(\omega) = -\frac{d\theta(\omega)}{d\omega}$

\subsection*{Step 2: Substitute Given Phase}
Given $\theta(\omega) = -3\omega$:
$\tau_p(\omega) = -\frac{-3\omega}{\omega} = 3\ \text{seconds}$
$\tau_g(\omega) = -\frac{d(-3\omega)}{d\omega} = 3\ \text{seconds}$
Since both are equal and constant, the system exhibits linear phase (perfect delay without phase distortion).

Correct Answer: C
\end{solution}

\clearpage

\subsection{Matched Filter SNR Maximization}

A matched filter is a specialized linear filter designed to maximize the:

\begin{choices}
    \choice Bandwidth of the incoming communication channel
    \choice Flatness of the filter magnitude response
    \choice Phase linearity of the output signal
    \correctchoice Signal-to-noise ratio (SNR) at a specific sampling instant in the presence of noise
\end{choices}

\begin{solution}

\subsection*{Step 1: Define Matched Filter Goal}
In communications and radar, the matched filter's impulse response is a time-reversed and delayed version of the target signal: $h(t) = s^*(T - t)$. This shape maximizes the peak output signal-to-noise ratio in the presence of additive white Gaussian noise.

Correct Answer: D
\end{solution}

\clearpage

\subsection{Energy Spectral Density vs Power Spectral Density}

For a power signal, the Power Spectral Density (PSD) is the Fourier transform of the signal's:

\begin{choices}
    \choice Cross-correlation function $R_{xy}(\tau)$
    \correctchoice Autocorrelation function $R_{xx}(\tau)$
    \choice Impulse response $h(t)$
    \choice Time-domain envelope
\end{choices}

\begin{solution}

\subsection*{Step 1: State Wiener-Khinchin Theorem}
The Wiener-Khinchin theorem states that the Power Spectral Density $S_{xx}(f)$ of a wide-sense stationary random process (or a power signal) is the Fourier transform of its temporal autocorrelation function $R_{xx}(\tau)$:
$S_{xx}(f) = \mathcal{F}\{R_{xx}(\tau)\} = \int_{-\infty}^{\infty} R_{xx}(\tau) e^{-j 2\pi f \tau}\, d\tau$

Correct Answer: B
\end{solution}

\clearpage

\subsection{Autocorrelation Peak}

The autocorrelation function $R_{xx}(\tau)$ of any signal $x(t)$ always achieves its absolute maximum value at:

\begin{choices}
    \choice $\tau = \infty$
    \choice $\tau$ equal to the fundamental period
    \choice The first zero-crossing of the signal
    \correctchoice $\tau = 0$
\end{choices}

\begin{solution}

\subsection*{Step 1: Prove Autocorrelation Maximum}
By the Cauchy-Schwarz inequality:
$|R_{xx}(\tau)| = \left| \int x(t) x(t-\tau)\, dt \right| \le \int x^2(t)\, dt = R_{xx}(0)$
Thus, $|R_{xx}(\tau)| \le R_{xx}(0)$ for all $\tau$, meaning the peak always occurs at zero lag (which represents the signal's total energy or power).

Correct Answer: D
\end{solution}

\clearpage

\subsection{Fourier Transform of Exponential Pulse}

What is the continuous-time Fourier transform of the signal $x(t) = e^{-3t} u(t)$?

\begin{choices}
    \choice $\frac{1}{3 - j\omega}$
    \choice $\frac{1}{9 + \omega^2}$
    \correctchoice $\frac{1}{3 + j\omega}$
    \choice $\frac{3}{3 + j\omega}$
\end{choices}

\begin{solution}

\subsection*{Step 1: Write Fourier Transform Integral}
$X(j\omega) = \int_{-\infty}^{\infty} x(t) e^{-j\omega t}\, dt = \int_{0}^{\infty} e^{-3t} e^{-j\omega t}\, dt$

\subsection*{Step 2: Evaluate Integral}
$X(j\omega) = \int_{0}^{\infty} e^{-(3 + j\omega)t}\, dt = \left[ -\frac{1}{3 + j\omega} e^{-(3 + j\omega)t} \right]_{0}^{\infty} = 0 - \left( -\frac{1}{3 + j\omega} \right) = \frac{1}{3 + j\omega}$

Correct Answer: C
\end{solution}

\clearpage

\subsection{Duality Property of Fourier Transform}

The duality property of the continuous Fourier transform states that if $\mathcal{F}\{x(t)\} = X(f)$, then what is the Fourier transform of the frequency-domain shape in the time domain, $\mathcal{F}\{X(t)\}$?

\begin{choices}
    \choice $x(f)$
    \choice $X(-f)$
    \choice $-x(f)$
    \correctchoice $x(-f)$
\end{choices}

\begin{solution}

\subsection*{Step 1: State Duality Property}
The symmetry (duality) of the forward and inverse Fourier transform integrals implies:
$\mathcal{F}\{X(t)\} = x(-f)$
or in terms of angular frequency $\omega$:
$\mathcal{F}\{X(t)\} = 2\pi x(-\omega)$

Correct Answer: D
\end{solution}

\clearpage

\subsection{Fourier Transform of Rectangular Pulse}

What is the Fourier transform $X(f)$ of a rectangular pulse of width $T$ and amplitude $A$, centered at the origin ($t=0$)?

\begin{choices}
    \correctchoice $A T \text{sinc}(f T)$
    \choice $A \text{sinc}(f T)$
    \choice $A T \text{sinc}^2(f T)$
    \choice $\frac{A}{j 2\pi f}$
\end{choices}

\begin{solution}

\subsection*{Step 1: Write Integral}
$X(f) = \int_{-T/2}^{T/2} A e^{-j 2\pi f t}\, dt$

\subsection*{Step 2: Evaluate Integral}
$X(f) = A \left[ \frac{e^{-j 2\pi f t}}{-j 2\pi f} \right]_{-T/2}^{T/2} = A \frac{e^{j \pi f T} - e^{-j \pi f T}}{j 2\pi f}$
Using Euler's identity:
$X(f) = A \frac{\sin(\pi f T)}{\pi f} = A T \frac{\sin(\pi f T)}{\pi f T} = A T \text{sinc}(f T)$

Correct Answer: A
\end{solution}

\clearpage

\subsection{Fourier Transform of Unit Impulse}

What is the continuous Fourier transform of the unit impulse function $x(t) = \delta(t)$?

\begin{choices}
    \choice $\delta(f)$
    \correctchoice 1 (constant for all frequencies)
    \choice $0$
    \choice $j 2\pi f$
\end{choices}

\begin{solution}

\subsection*{Step 1: Apply Delta Integral Property}
$X(f) = \int_{-\infty}^{\infty} \delta(t) e^{-j 2\pi f t}\, dt$

\subsection*{Step 2: Sift at $t = 0$}
The sifting property evaluates the integrand at $t = 0$:
$X(f) = e^{-j 2\pi f (0)} = e^0 = 1$

Correct Answer: B
\end{solution}

\clearpage

\subsection{Frequency Shifting (Modulation)}

Multiplying a time-domain signal by a complex exponential $e^{j 2\pi f_0 t}$ corresponds in the frequency domain to:

\begin{choices}
    \choice Shifting the spectrum by $-f_0$, $X(f + f_0)$
    \choice Scaling the frequency axis, $X(f / f_0)$
    \correctchoice Shifting the spectrum by $f_0$, $X(f - f_0)$
    \choice Convolving the spectrum with $f_0$
\end{choices}

\begin{solution}

\subsection*{Step 1: Apply Frequency Shifting Property}
The modulation property states:
$\mathcal{F}\{x(t) e^{j 2\pi f_0 t}\} = \int_{-\infty}^{\infty} x(t) e^{j 2\pi f_0 t} e^{-j 2\pi f t}\, dt$
$= \int_{-\infty}^{\infty} x(t) e^{-j 2\pi (f - f_0) t}\, dt = X(f - f_0)$

Correct Answer: C
\end{solution}

\clearpage

\subsection{Parseval's Relation for Fourier Transform}

Parseval's theorem for the continuous Fourier transform states that the total energy of a signal is conserved. This is mathematically written as:

\begin{choices}
    \choice $\int_{-\infty}^{\infty} |x(t)|^2\, dt = \int_{-\infty}^{\infty} |X(\omega)|^2\, d\omega$
    \choice $\int_{-\infty}^{\infty} x(t)\, dt = \int_{-\infty}^{\infty} X(f)\, df$
    \choice $\int_{-\infty}^{\infty} |x(t)|^2\, dt = \frac{1}{2\pi} \int_{-\infty}^{\infty} |X(f)|^2\, df$
    \correctchoice $\int_{-\infty}^{\infty} |x(t)|^2\, dt = \int_{-\infty}^{\infty} |X(f)|^2\, df$
\end{choices}

\begin{solution}

\subsection*{Step 1: State Energy Conservation}
Parseval's identity states that the signal energy in the time domain is equal to the signal energy in the frequency domain:
$\int_{-\infty}^{\infty} |x(t)|^2\, dt = \int_{-\infty}^{\infty} |X(f)|^2\, df = \frac{1}{2\pi} \int_{-\infty}^{\infty} |X(j\omega)|^2\, d\omega$

Correct Answer: D
\end{solution}

\clearpage

\subsection{Fourier Transform of DC Signal}

What is the continuous Fourier transform of a constant DC voltage $x(t) = V_0$?

\begin{choices}
    \correctchoice $V_0 \delta(f)$
    \choice $V_0$
    \choice $0$
    \choice $\frac{V_0}{j 2\pi f}$
\end{choices}

\begin{solution}

\subsection*{Step 1: Apply Duality to Impulse}
We know that $\mathcal{F}\{\delta(t)\} = 1$. By the duality property:
$\mathcal{F}\{1\} = \delta(-f) = \delta(f)$

\subsection*{Step 2: Scale by $V_0$}
Since the Fourier transform is a linear operator:
$\mathcal{F}\{V_0\} = V_0 \delta(f)$

Correct Answer: A
\end{solution}

\clearpage

\subsection{Fourier Transform Time Scaling Property}

If the Fourier transform of $x(t)$ is $X(f)$, what is the Fourier transform of the time-scaled signal $y(t) = x(a t)$ for $a 

e 0$?

\begin{choices}
    \choice $|a| X(a f)$
    \correctchoice $\frac{1}{|a|} X\left(\frac{f}{a}\right)$
    \choice $\frac{1}{a} X(a f)$
    \choice $X\left(\frac{f}{a}\right)$
\end{choices}

\begin{solution}

\subsection*{Step 1: Apply Time Scaling Integral}
$\mathcal{F}\{x(at)\} = \int_{-\infty}^{\infty} x(at) e^{-j 2\pi f t}\, dt$

\subsection*{Step 2: Evaluate by Substitution}
Let $u = at \implies dt = du/a$. If $a > 0$, limits are unchanged. If $a < 0$, limits are swapped, multiplying by $-1$. Combining both cases using absolute value:
$\mathcal{F}\{x(at)\} = \frac{1}{|a|} \int_{-\infty}^{\infty} x(u) e^{-j 2\pi (f/a) u}\, du = \frac{1}{|a|} X\left(\frac{f}{a}\right)$

Correct Answer: B
\end{solution}

\clearpage

\subsection{Fourier Transform Differentiation Property}

If $X(f)$ is the Fourier transform of $x(t)$, what is the Fourier transform of the derivative $\frac{dx(t)}{dt}$?

\begin{choices}
    \choice $\frac{X(f)}{j 2\pi f}$
    \choice $j\omega X(f)$
    \correctchoice $j 2\pi f X(f)$
    \choice $-j 2\pi f X(f)$
\end{choices}

\begin{solution}

\subsection*{Step 1: Apply Derivative Property}
Differentiating in the time domain corresponds to multiplying by $j 2\pi f$ (or $j\omega$) in the frequency domain:
$\mathcal{F}\left\{\frac{dx(t)}{dt}\right\} = j 2\pi f X(f)$

Correct Answer: C
\end{solution}

\clearpage

\subsection{Cross-Correlation}

A massive radar station blasts a radio pulse at an airplane. The pulse bounces off the plane and returns as a tiny, noisy echo. To physically calculate exactly how many microseconds it took for the echo to return, the computer mathematically performs 'Cross-Correlation' between the original pulse and the noisy sky. What does this math do?

\begin{choices}
    \choice It deletes the signal
    \choice It multiplies them by zero
    \correctchoice It mathematically physically slides the clean original pulse across the noisy radar data; at the exact millisecond the two shapes perfectly 'lock' together, the mathematical equation violently spikes, instantly revealing the precise time delay
    \choice It divides the frequencies
\end{choices}

\begin{solution}

\subsection*{Step 1: The Needle in the Haystack}
The returning echo is so weak it is buried under $10\text{ dB}$ of random atmospheric static. The human eye cannot see it.

\subsection*{Step 2: The Pattern Matcher}
Cross-Correlation is mathematically just 'Pattern Matching'. It slides a template of the original pulse across the static. When it hits random static, the multiplications cancel out to near zero. But the exact nanosecond the template physically overlays perfectly on top of the hidden echo, all the $1$s line up with $1$s, and the math explodes into a massive $+100$ spike. The radar instantly knows the plane is $14.2\text{ miles}$ away.

Correct Answer: C
\end{solution}

\clearpage

\subsection{Phase Delay vs Group Delay}

When a complex signal (like a human voice) travels through a physical electronic filter, the filter mathematically delays the wave. If the filter has 'Linear Phase', what mathematically and physically happens to the human voice?

\begin{choices}
    \choice It turns into an alien voice
    \choice It echoes repeatedly
    \correctchoice The voice sounds physically flawless; Linear Phase mathematically guarantees that ALL frequencies (bass and treble) are delayed by the exact same physical amount of milliseconds (Constant Group Delay), preventing the wave from smearing
    \choice It mutes the bass
\end{choices}

\begin{solution}

\subsection*{Step 1: The Smear (Non-Linear Phase)}
A human voice is made of a $100\text{ Hz}$ bass wave and a $2000\text{ Hz}$ treble wave hitting the microphone at the exact same millisecond. If a bad filter delays the bass by $10\text{ ms}$, but the treble by $50\text{ ms}$, the voice physically stretches out like a slinky. It sounds unrecognizable.

\subsection*{Step 2: The Linear Perfection}
If the Phase math is perfectly 'Linear' (a straight line graph), then Phase / Frequency = Constant. This means every single tone is delayed by exactly $20\text{ ms}$. The entire voice wave arrives late, but it arrives in mathematically perfect unison. FIR digital filters are famous for being able to physically achieve absolute mathematically perfect Linear Phase.

Correct Answer: C
\end{solution}

\clearpage

\subsection{Efficiency}

Transformer efficiency is generally highest at:

\begin{choices}
    \choice No load
    \correctchoice Full load or near full load
    \choice Quarter load
    \choice Twice the rated load
\end{choices}

\begin{solution}

\subsection*{Step 1: Condition}
Occurs when copper losses equal core losses.

Correct Answer: B
\end{solution}

\clearpage

\subsection{Core Losses}

Transformer core losses consist of:

\begin{choices}
    \choice Dielectric losses
    \choice Copper losses
    \choice Stray load losses
    \correctchoice Hysteresis and Eddy current losses
\end{choices}

\begin{solution}

\subsection*{Step 1: Nature}
Independent of load current (constant losses).

Correct Answer: D
\end{solution}

\clearpage

\subsection{Dot Convention}

The dot convention in transformers is used to indicate:

\begin{choices}
    \choice The center of the winding
    \choice The primary side
    \choice Temperature limit
    \correctchoice Relative polarity of voltages
\end{choices}

\begin{solution}

\subsection*{Step 1: Logic}
Determines whether voltages are in phase or 180° out of phase.

Correct Answer: D
\end{solution}

\clearpage

\subsection{Real Transformer Voltage Regulation Sizing}

A transformer is connected to a primary supply. The secondary terminal voltage is measured to be $V_{noload} = 240\text{ V}$ when open-circuited. Under full-load conditions, the secondary terminal voltage drops to $V_{fullload} = 230\text{ V}$. Calculate the percentage voltage regulation ($VR\%$).

\begin{choices}
    \correctchoice 4.35%
    \choice 4.17%
    \choice 9.52%
    \choice 5.00%
\end{choices}

\begin{solution}

\subsection*{Step 1: Identify Voltage Regulation Formula}
$VR\% = \frac{V_{noload} - V_{fullload}}{V_{fullload}} \times 100\%$

\subsection*{Step 2: Calculate Percentage}
$VR\% = \frac{240 - 230}{230} \times 100\% = \frac{10}{230} \times 100\% = 4.35\%$.

Correct Answer: A
\end{solution}

\clearpage

\subsection{Equivalent Core Reactance}

In a $100\text{ kVA}$, $2400/240\text{ V}$ step-down transformer, the open-circuit test is performed on the low-voltage side. The test readings are $V_{oc} = 240\text{ V}$, $I_{oc} = 2\text{ A}$, and $P_{oc} = 150\text{ W}$. Calculate the magnetizing shunt reactance ($X_m$) referred to the low-voltage side.

\begin{choices}
    \correctchoice 126.5 \Omega
    \choice 120.0 \Omega
    \choice 384.0 \Omega
    \choice 72.4 \Omega
\end{choices}

\begin{solution}

\subsection*{Step 1: Calculate Core Loss Branch Admittance}
No-load phase angle: $pf_{oc} = \cos(\theta_{oc}) = \frac{P_{oc}}{V_{oc} I_{oc}} = \frac{150}{240 \times 2} = \frac{150}{480} = 0.3125$.
$\theta_{oc} = \cos^{-1}(0.3125) = -71.79^circ$.

\subsection*{Step 2: Calculate Magnetizing Current}
No-load current in polar form: $\mathbf{I}_{oc} = 2.0 \angle -71.79^circ\text{ A}$.
Magnetizing reactive current component: $I_m = I_{oc} \sin(\theta_{oc}) = 2.0 \sin(71.79^circ) = 2.0 \times 0.9499 = 1.90\text{ A}$.

\subsection*{Step 3: Calculate Magnetizing Reactance}
$X_m = \frac{V_{oc}}{I_m} = \frac{240}{1.90} = 126.3\text{ \Omega} \approx 126.5\text{ \Omega}$.

Correct Answer: A
\end{solution}

\clearpage

\subsection{Equivalent Winding Impedance}

A short-circuit test is performed on the high-voltage side of a transformer. The test readings are $V_{sc} = 50\text{ V}$, $I_{sc} = 10\text{ A}$, and $P_{sc} = 300\text{ W}$. Calculate the equivalent series resistance ($R_{eq}$) of the transformer winding referred to the high-voltage side.

\begin{choices}
    \correctchoice 3.0 \Omega
    \choice 5.0 \Omega
    \choice 4.0 \Omega
    \choice 2.0 \Omega
\end{choices}

\begin{solution}

\subsection*{Step 1: Identify Short-circuit Resistance Relation}
The short-circuit power reading represents the copper losses in the windings: $P_{sc} = I_{sc}^2 R_{eq}$.

\subsection*{Step 2: Solve for Req}
$R_{eq} = \frac{P_{sc}}{I_{sc}^2} = \frac{300}{10^2} = \frac{300}{100} = 3.0\text{ \Omega}$.

Correct Answer: A
\end{solution}

\clearpage

\subsection{Maximum Efficiency Load Condition}

A transformer has core losses of $200\text{ W}$ and full-load copper losses of $800\text{ W}$. At what fraction ($x$) of full-load rated kVA will the transformer operate at its maximum possible efficiency?

\begin{choices}
    \correctchoice 0.50 (50% load)
    \choice 0.25 (25% load)
    \choice 0.75 (75% load)
    \choice 1.00 (100% load)
\end{choices}

\begin{solution}

\subsection*{Step 1: State Maximum Efficiency Condition}
Maximum efficiency in a transformer occurs when the variable copper losses ($x^2 P_{cu,fl}$) equal the constant core losses ($P_{core}$):
$x^2 P_{cu,fl} = P_{core}$.

\subsection*{Step 2: Solve for x}
$x = \sqrt{\frac{P_{core}}{P_{cu,fl}}} = \sqrt{\frac{200}{800}} = \sqrt{0.25} = 0.50$.

Correct Answer: A
\end{solution}

\clearpage

\subsection{Equivalent Winding Leakage Reactance}

For the short-circuit test in a previous question ($V_{sc} = 50\text{ V}$, $I_{sc} = 10\text{ A}$, $R_{eq} = 3\text{ \Omega}$), calculate the equivalent series leakage reactance ($X_{eq}$) of the windings referred to the test side.

\begin{choices}
    \correctchoice 4.0 \Omega
    \choice 5.0 \Omega
    \choice 3.0 \Omega
    \choice 2.5 \Omega
\end{choices}

\begin{solution}

\subsection*{Step 1: Calculate Equivalent Impedance magnitude}
$Z_{eq} = \frac{V_{sc}}{I_{sc}} = \frac{50}{10} = 5.0\text{ \Omega}$.

\subsection*{Step 2: Solve for Xeq}
$X_{eq} = \sqrt{Z_{eq}^2 - R_{eq}^2} = \sqrt{5.0^2 - 3.0^2} = \sqrt{25 - 9} = \sqrt{16} = 4.0\text{ \Omega}$.

Correct Answer: A
\end{solution}

\clearpage

\subsection{Real Transformer Core Loss Components}

Which physical mechanism accounts for the core losses inside a real transformer working under steady-state alternating magnetic fields?

\begin{choices}
    \correctchoice Both magnetic hysteresis loops and induced eddy current heating
    \choice Ohmic resistance heating in the copper conductors alone
    \choice Mechanical vibration in the mounting bracket structure
    \choice Dielectric breakdown in the winding isolation paper
\end{choices}

\begin{solution}

\subsection*{Step 1: Analyze core loss mechanisms}
Core losses in ferromagnetic laminations are composed of:
1. Hysteresis loss: Energy required to continually align magnetic domains under alternating AC fields.
2. Eddy current loss: Ohmic $I^2R$ heating losses due to circulating currents induced in the conductive iron core by alternating flux.

Correct Answer: A
\end{solution}

\clearpage

\subsection{Three-Phase Transformer Phase Displacement}

A three-phase transformer bank is connected in Wye-Delta ($Y-\Delta$). What is the phase angle displacement introduced between the primary line-to-line voltages and secondary line-to-line voltages?

\begin{choices}
    \correctchoice 30 degrees phase shift
    \choice 0 degrees phase shift
    \choice 120 degrees phase shift
    \choice 90 degrees phase shift
\end{choices}

\begin{solution}

\subsection*{Step 1: Understand Y-Delta Displacement}
By IEEE standard, three-phase transformer connections involving mismatched types (Wye-Delta or Delta-Wye) introduce a fundamental phase displacement of $30^\circ$ between the corresponding line voltages.

Correct Answer: A
\end{solution}

\clearpage

\subsection{Transformer Efficiency}

A 50 kVA transformer has iron losses of 500 W and full-load copper losses of 800 W. At what fraction of full load does maximum efficiency occur?

\begin{choices}
    \correctchoice 79.1%
    \choice 62.5%
    \choice 100.0%
    \choice 50.0%
\end{choices}

\begin{solution}

\subsection*{Step 1: Max Efficiency Condition}
Maximum efficiency occurs when variable copper losses equal constant iron losses. $x^2 P_{cu} = P_{core} \Rightarrow x = \sqrt{\dfrac{P_{core}}{P_{cu}}} = \sqrt{\dfrac{500}{800}} = \sqrt{0.625} \approx 0.791$ or 79.1%.

Correct Answer: A
\end{solution}

\clearpage

\subsection{SNR Calculation}

If the signal power is 1000 W and the noise power is 1 W, the SNR in dB is:

\begin{choices}
    \choice 60 dB
    \choice 1000 dB
    \correctchoice 30 dB
    \choice 10 dB
\end{choices}

\begin{solution}

\subsection*{Step 1: Formula}
$SNR_{dB} = 10 log_{10}(P_s / P_n) = 10 log_{10}(1000) = 30$.

Correct Answer: C
\end{solution}

\clearpage

\subsection{Thermal Noise}

Thermal noise is also known as:

\begin{choices}
    \choice Pink noise
    \choice Impulse noise
    \correctchoice White noise
    \choice Flicker noise
\end{choices}

\begin{solution}

\subsection*{Step 1: Nature}
Has a constant power spectral density over all frequencies.

Correct Answer: C
\end{solution}

\clearpage

\subsection{SNR in Decibels Calculation}

A receiver detects a signal power of $P_{sig} = 10\text{ mW}$ and a noise power of $P_{noise} = 2.0\text{ \mu W}$. Calculate the Signal-to-Noise Ratio ($SNR$) in decibels (dB).

\begin{choices}
    \correctchoice 37.0 dB
    \choice 5000 dB
    \choice 23.0 dB
    \choice 30.0 dB
\end{choices}

\begin{solution}

\subsection*{Step 1: Calculate Linear SNR Ratio}
$SNR_{linear} = \frac{P_{sig}}{P_{noise}} = \frac{10 \times 10^{-3}\text{ W}}{2.0 \times 10^{-6}\text{ W}} = 5000$.

\subsection*{Step 2: Convert to Decibels}
$SNR_{dB} = 10 \log_{10}(SNR_{linear}) = 10 \log_{10}(5000) \approx 10 \times 3.69897 = 36.99\text{ dB}$.

Correct Answer: A
\end{solution}

\clearpage

\subsection{Thermal Noise Power Sizing}

Calculate the thermal noise power ($P_N$) in Watts generated by a resistor at room temperature ($T = 290\text{ K}$) over an audio bandwidth of $B = 20\text{ kHz}$. (Boltzmann's constant $k = 1.38 \times 10^{-23}\text{ J/K}$)

\begin{choices}
    \correctchoice 8.00 \times 10^{-17} W
    \choice 4.00 \times 10^{-17} W
    \choice 8.00 \times 10^{-14} W
    \choice 1.16 \times 10^{-16} W
\end{choices}

\begin{solution}

\subsection*{Step 1: Identify Thermal Noise Power Formula}
$P_N = k T B$

\subsection*{Step 2: Calculate Value}
$P_N = (1.38 \times 10^{-23}\text{ J/K}) \times 290\text{ K} \times (20 \times 10^3\text{ Hz})$
$P_N = 4.002 \times 10^{-21} \times 20000 = 8.004 \times 10^{-17}\text{ Watts}$.

Correct Answer: A
\end{solution}

\clearpage

\subsection{Noise Figure Sizing}

An amplifier has an input Signal-to-Noise Ratio of $SNR_{in} = 40\text{ dB}$ and an output Signal-to-Noise Ratio of $SNR_{out} = 37\text{ dB}$. Calculate the Noise Figure ($NF$) of the amplifier in dB.

\begin{choices}
    \correctchoice 3 dB
    \choice 77 dB
    \choice 1.08 dB
    \choice 6 dB
\end{choices}

\begin{solution}

\subsection*{Step 1: Identify Noise Figure in Decibels}
$NF_{dB} = SNR_{in,dB} - SNR_{out,dB}$

\subsection*{Step 2: Calculate Value}
$NF_{dB} = 40 - 37 = 3\text{ dB}$ (which represents a linear Noise Factor of $F = 2.0$).

Correct Answer: A
\end{solution}

\clearpage

\subsection{Noise Factor Cascaded Stages (Friis Formula)}

A receiver contains two cascaded amplifier stages. Stage 1 has a noise factor of $F_1 = 2.0$ and a gain of $G_1 = 10.0$. Stage 2 has a noise factor of $F_2 = 4.0$. Calculate the total receiver noise factor ($F_{total}$) using the Friis Formula.

\begin{choices}
    \correctchoice 2.3
    \choice 6.0
    \choice 2.4
    \choice 3.0
\end{choices}

\begin{solution}

\subsection*{Step 1: Identify Friis Formula for Cascaded Noise}
$F_{total} = F_1 + \frac{F_2 - 1}{G_1}$

\subsection*{Step 2: Substitute and Calculate Factor}
$F_{total} = 2.0 + \frac{4.0 - 1}{10.0} = 2.0 + \frac{3}{10.0} = 2.0 + 0.3 = 2.3$.

Correct Answer: A
\end{solution}

\clearpage

\subsection{Effective Noise Temperature}

An amplifier stage has a noise factor of $F = 2.0$. What is the equivalent effective noise temperature ($T_e$) of the amplifier, referencing room temperature at $T_0 = 290\text{ K}$?

\begin{choices}
    \correctchoice 290 K
    \choice 580 K
    \choice 0 K
    \choice 145 K
\end{choices}

\begin{solution}

\subsection*{Step 1: Identify Noise Temperature Formula}
$T_e = T_0 (F - 1)$

\subsection*{Step 2: Calculate Temperature}
$T_e = 290\text{ K} \times (2.0 - 1) = 290\text{ K}$.

Correct Answer: A
\end{solution}

\clearpage

\subsection{Additive White Gaussian Noise (AWGN) Definition}

What is the physical meaning of the 'white' and 'Gaussian' descriptors in an Additive White Gaussian Noise (AWGN) channel model?

\begin{choices}
    \correctchoice White indicates flat power spectral density; Gaussian indicates normal probability distribution of voltage values
    \choice White indicates zero thermal power; Gaussian indicates purely reactive noise impedance
    \choice White indicates polar phase angles; Gaussian indicates periodic ripple components
    \choice Both terms indicate a single constant noise frequency component
\end{choices}

\begin{solution}

\subsection*{Step 1: Deconstruct AWGN Channel Terms}
- **White**: The noise has uniform (flat) power spectral density across all frequencies, analogous to white light containing all visible frequencies.
- **Gaussian**: The instantaneous amplitude of the noise voltage follows a Gaussian (normal) probability density function with zero mean.
- **Additive**: The noise voltage is simply added to the transmitted signal.

Correct Answer: A
\end{solution}

\clearpage

\subsection{Resistor Noise Equivalent Model}

Under the Thevenin equivalent model of a real noisy resistor $R$, what is the RMS value of the open-circuit thermal noise voltage ($v_n$) generated by the resistor at temperature $T$ over bandwidth $B$?

\begin{choices}
    \correctchoice \sqrt{4 k T R B}
    \choice 4 k T R B
    \choice \sqrt{k T B}
    \choice \sqrt{\frac{4 k T B}{R}}
\end{choices}

\begin{solution}

\subsection*{Step 1: State Johnson-Nyquist Noise Formula}
The thermal noise generated by a resistor is modeled as a noise voltage source in series with an ideal noiseless resistor: $v_{n}^2 = 4 k T R B \implies v_n = \sqrt{4 k T R B}$.

Correct Answer: A
\end{solution}

\clearpage

\subsection{Principal Stresses on Free Surface}

On a plane where shear stress is zero, the normal stresses acting on that plane are called:

\begin{choices}
    \correctchoice Principal stresses
    \choice Shear stresses
    \choice Yield stresses
    \choice Deviatoric stresses
\end{choices}

\begin{solution}

\subsection*{Step 1: Principal Stresses}
By definition, principal planes are planes of zero shear stress, and the normal stresses on these planes are the maximum and minimum principal stresses.

Correct Answer: A
\end{solution}

\clearpage

